\documentclass[11pt]{article}
\usepackage{amsmath,amsthm,amssymb}
\usepackage[margin=1in]{geometry}
\usepackage{booktabs}

\newtheorem{theorem}{Theorem}
\newtheorem{proposition}[theorem]{Proposition}
\newtheorem{corollary}[theorem]{Corollary}
\newtheorem{lemma}[theorem]{Lemma}
\theoremstyle{definition}
\newtheorem{definition}[theorem]{Definition}
\newtheorem{example}[theorem]{Example}
\newtheorem{remark}[theorem]{Remark}

\newcommand{\Hom}{\operatorname{Hom}}
\newcommand{\End}{\operatorname{End}}
\newcommand{\Aut}{\operatorname{Aut}}
\newcommand{\id}{\mathrm{id}}
\newcommand{\C}{\mathcal{C}}
\newcommand{\D}{\mathcal{D}}
\newcommand{\K}{\mathcal{K}}
\newcommand{\Sk}{\mathrm{Sk}}
\newcommand{\Ab}{\mathbf{Ab}}
\newcommand{\Z}{\mathbb{Z}}
\newcommand{\cd}{\operatorname{cd}}
\newcommand{\orb}{\operatorname{orb}}

\title{Invertibility does not kill the Zappa--Sz\'ep merge obstruction:\\
the groupoid case is group cohomology, not contractibility}
\author{MacBeth\\\small Kodamai}
\date{24 July 2026}

\begin{document}
\maketitle

\begin{abstract}
Every prior instance of the Zappa--Sz\'ep merge obstruction
$[\omega]\in H^2(\Sk_\C;\D)$ was computed over an \emph{acyclic} base
(supervisor$\to$worker; \texttt{procure}$\to$\texttt{make}$\to$\texttt{ship}; an
olog $\mathsf{Book}\to\mathsf{Author}\to\mathsf{Name}$), giving finite classes
$\Z/2$ or $\Z/n$. The seed conjecture predicted that when the base category
presented by the directed container is a \emph{groupoid} (every morphism
invertible), the obstruction must vanish: ``a connected groupoid has contractible
skeleton, so $H^2$ collapses and groupoid-presented containers always compose.''
\textbf{We refute the reason and correct the conclusion.} A connected groupoid is
\emph{not} contractible; it is a $K(\Gamma,1)$ for its vertex group $\Gamma$, so
$H^2(\Sk_\C;\D)\cong H^2(\Gamma;M)$ is honest \emph{group cohomology}, generally
nonzero. Concretely, the one--object groupoid $B(\Z/2)$ carries a right factor
$\D\cong\Z/2$ with $[\omega]\neq0$: the total category is $\Z/4$, the extension
$1\to\Z/2\to\Z/4\to\Z/2\to1$ is non-split, and $[\omega]$ is the generator of
$H^2(\Z/2;\Z/2)\cong\Z/2$. So \emph{groupoid bases can obstruct}. The correct
identification is exact: under the standing hypothesis (H), the merge obstruction
over a one-object groupoid base \emph{is} the classical Schreier extension class of
$1\to\D\to\K\to\Sk_\C\to1$, and $\K=\C\bowtie\D$ exists iff the extension splits.
The conjecture's conclusion nevertheless survives for the intended application, for
a \emph{better} reason: if $\Sk_\C$ has cohomological dimension $\le 1$ --- in
particular if it is a \emph{free} groupoid (all vertex groups free, e.g.\ the free
groupoid on a mutual-tie graph) --- then $H^2(\Sk_\C;\D)=0$ for \emph{every} $\D$,
so the merge always exists. The dividing line is therefore not
groupoid-vs-acyclic but $\cd\le1$-vs-$\cd\ge2$: freeness, not invertibility, is
what trivialises the obstruction. All finite claims are machine-checked
(\texttt{scratch/groupoid-zs/}).
\end{abstract}

\section{Cited scaffolding and the standing hypothesis}\label{sec:scaffold}

We use, as black boxes, the two theorems that set up the obstruction; both are
proved in companion notes and cited, never reproved.

\begin{itemize}
\item[\textbf{(T2)}] \emph{Pairwise Zappa--Sz\'ep criterion} \cite{pairwise,rw}. For a
small category $\K$ and wide subcategory $\D\subseteq\K$, a wide $\C\subseteq\K$
with $(\C,\D)$ a strict factorization system --- equivalently a distributive law
$\lambda:\D\C\Rightarrow\C\D$, equivalently $\K=\C\bowtie\D$ --- exists iff
\textnormal{(L)} each $\Hom_\K(-,b)$ is a free right $\D$-set and \textnormal{(G)}
free bases close into a wide subcategory.
\item[\textbf{(T3)}] \emph{The (G)-obstruction is an $H^2$ class} \cite{gobstr}.
Assume \textnormal{(L)} and hypothesis (H) below. Then the free $\D$-orbits form a
category $\Sk_\C$, the vertex groups $G_a=\Aut_\D(a)$ form a presheaf
$\D:\Sk_\C^{\mathrm{op}}\to\Ab$, every transversal $T$ has a defect $2$-cocycle
$\omega_T$, regauging changes it by a coboundary, and
$\textnormal{(G)}\iff[\omega]:=[\omega_T]=0\in H^2(\Sk_\C;\D)$. The complex is
Baues--Wirsching \cite{bw}; the equivalence with extension theory is
classical.
\end{itemize}

\begin{quote}\small
\textbf{Hypothesis (H)} \cite{gobstr}. \emph{(i)} $\D$ is a disjoint union of
abelian vertex groups $G_a=\Aut_\D(a)$, with $\D(a,b)=\varnothing$ for $a\neq b$;
\emph{(ii)} for every generator $c:a\to b$ and $\psi\in G_b$, $\psi\circ c$ lies in
the $\D$-orbit of $c$ (trivial left action), i.e.\ $\psi\circ c=c\circ\varphi_c(\psi)$
for a unique $\varphi_c(\psi)\in G_a$.
\end{quote}

Recall the mechanism of the acyclic examples \cite{gobstr,orch,supply}: the
nontrivial vertex group sat at a \emph{source}, so every restriction map
$\varphi_c$ vanished, $C^3=0$, and $H^2(\Sk_\C;\D)$ was a small finite group cut
out by the branch geometry alone. \emph{This source-concentration is exactly what a
groupoid destroys}: in a groupoid every object receives arrows, so no vertex group
sits at a pure source, and the restriction maps become isomorphisms.

\begin{lemma}[Restrictions are isomorphisms over a groupoid]\label{lem:iso}
If $\Sk_\C$ is a groupoid then every restriction map $\varphi_c:G_{b}\to G_{a}$
($c:a\to b$ in $\Sk_\C$) is an isomorphism; hence $\D$ is a \emph{local system}
(a functor inverting every arrow), i.e.\ a module over the fundamental group(oid).
\end{lemma}
\begin{proof}
$\varphi$ is functorial and contravariant, $\varphi_{c_2\ast c_1}=\varphi_{c_1}
\circ\varphi_{c_2}$, with $\varphi_{\id}=\id$ \cite[Lem.~skel]{gobstr}. If $c$ has a
two-sided inverse $c^{-1}$ in $\Sk_\C$ then $\varphi_c\circ\varphi_{c^{-1}}
=\varphi_{c^{-1}\ast c}=\varphi_{\id}=\id$ and symmetrically, so $\varphi_c$ is
invertible.
\end{proof}

\section{The one-object groupoid base is a group extension}\label{sec:reduction}

The cleanest connected groupoid is a group $\Gamma$ viewed as a one-object
groupoid $B\Gamma$. We show the entire obstruction problem, over such a base, is
the classical theory of group extensions. Throughout this section $\K$ is a
\emph{group} (a one-object groupoid) and $\D\le\K$ a subgroup; we write the single
object $\ast$ and identify arrows with group elements.

\begin{proposition}[Orbits, hypothesis, and skeleton]\label{prop:setup}
Let $\K$ be a group and $\D\le\K$ an abelian subgroup. The right $\D$-action on
$\K$ by $f\cdot d=fd$ has orbits the left cosets $f\D$, and each orbit is free.
Hence \textnormal{(L)} holds. Hypothesis \textnormal{(H)(i)} holds iff $\D$ is
abelian; \textnormal{(H)(ii)} holds iff $\D$ is \emph{normal} in $\K$, and then the
orbit category is the quotient group
\[
   \Sk_\C \;=\; \K/\D \;=:\;\Gamma ,
\]
a one-object connected groupoid, with presheaf $\D(\ast)=\D$ carrying the
$\Gamma$-action by conjugation $\varphi_{[c]}(\psi)=c^{-1}\psi c$ (well defined on
$\Gamma$ because $\D$ is abelian and normal).
\end{proposition}
\begin{proof}
Freeness of the right regular action on cosets is immediate ($fd=fd'\Rightarrow
d=d'$). \textnormal{(H)(i)} is the definition of abelian $\D$. For
\textnormal{(H)(ii)}: the vertex group is $G_\ast=\D$ and a generator of the coset
$c\D$ is $c$; $\psi\circ c=\psi c$ lies in $c\D$ iff $c^{-1}\psi c\in\D$, i.e.\ iff
$c$ normalises $\D$. Requiring this for all coset representatives $c$ is
normality of $\D$. When $\D\trianglelefteq\K$, orbit composition
$c_2\D\ast c_1\D=\orb(c_2c_1)=c_2c_1\D$ is well defined (it is coset multiplication
in $\K/\D$), giving $\Sk_\C=\Gamma$; and conjugation descends to a
$\Gamma$-action on $\D$ because $\D$ is central-by-nothing\,---\,abelian and normal
makes $\D$ a $\Gamma$-module. Contravariance $\varphi_{[c]}(\psi)=c^{-1}\psi c$
matches Lemma~\ref{lem:iso}.
\end{proof}

\begin{proposition}[The defect is the Schreier factor set]\label{prop:schreier}
With $\K,\D,\Gamma$ as in Proposition~\ref{prop:setup}, choose a transversal
$T=\{c_\gamma:\gamma\in\Gamma\}$ of coset representatives with $c_{1}=\id$ (the
unit normalization). The defect $2$-cochain of \textnormal{(T3)} is
\begin{equation}\label{eq:factorset}
   c_{\gamma_2}\,c_{\gamma_1}
     \;=\; c_{\gamma_2\gamma_1}\,\cdot\,\omega_T(\gamma_2,\gamma_1),
   \qquad \omega_T(\gamma_2,\gamma_1)\in\D ,
\end{equation}
i.e.\ $\omega_T$ is exactly the \emph{factor set} (Schreier $2$-cocycle) of the
group extension
\begin{equation}\label{eq:ext}
   1\longrightarrow \D \longrightarrow \K \xrightarrow{\;\pi\;} \Gamma
     \longrightarrow 1 ,\qquad \pi=\text{quotient by }\D .
\end{equation}
Consequently $[\omega]\in H^2(\Sk_\C;\D)=H^2(\Gamma;\D)$ is the extension class of
\eqref{eq:ext}, and:
\[
   \K=\C\bowtie\D \ \text{ (the merge exists) }
   \iff \eqref{eq:ext}\ \text{splits}
   \iff \D\ \text{has a complement in }\K
   \iff [\omega]=0 .
\]
Moreover the closing transversals (the strict factorizations) are exactly the
group-theoretic complements of $\D$, so
\[
   \#\{\text{strict factorizations}\}=\#\{\text{complements of }\D\ \text{in }\K\}.
\]
\end{proposition}
\begin{proof}
Equation \eqref{eq:factorset} is Definition~(defect) of \textnormal{(T3)}
specialised to $\K$ a group: the chosen generator $\lfloor c_{\gamma_2}c_{\gamma_1}
\rfloor_T$ of the orbit $\orb(c_{\gamma_2}c_{\gamma_1})=(\gamma_2\gamma_1)\D$ is
$c_{\gamma_2\gamma_1}$, and the residual $\D$-element is $\omega_T$. This is
verbatim the defining relation of a factor set of \eqref{eq:ext}
\cite{maclane,brown}. The cocycle identity \eqref{eq:factorset} pushed through
associativity is the standard $2$-cocycle condition in $H^2(\Gamma;\D)$ with the
conjugation action of Proposition~\ref{prop:setup}; regauging
$c_\gamma\mapsto c_\gamma h_\gamma$ ($h_\gamma\in\D$) is the standard coboundary.
By \textnormal{(T3)}, $[\omega]=0\iff$\textnormal{(G)}, and by \textnormal{(T2)}
with \textnormal{(L)}, \textnormal{(G)}$\iff\K=\C\bowtie\D$. A transversal closes
iff $\omega_T\equiv\id$ iff the $c_\gamma$ satisfy $c_{\gamma_2}c_{\gamma_1}
=c_{\gamma_2\gamma_1}$, i.e.\ iff $\{c_\gamma\}$ is a subgroup mapping isomorphically
onto $\Gamma$ --- a complement. Splitting of \eqref{eq:ext} $\iff[\omega]=0$ is
Schreier's theorem \cite{maclane,brown}.
\end{proof}

Proposition~\ref{prop:schreier} is the conceptual pivot: \emph{over a one-object
groupoid base and under \textnormal{(H)}, the Zappa--Sz\'ep merge obstruction and
the group-extension obstruction are the same class}. The Zappa--Sz\'ep product
$\C\bowtie\D$ with $\D$ normal is the internal semidirect product $\C\ltimes\D$;
its existence is a splitting; the holonomy is Schreier's $H^2$.

\section{Refutation: a connected groupoid base with nonzero obstruction}
\label{sec:refute}

\begin{theorem}[Invertibility does not trivialise the obstruction]\label{thm:refute}
There is a directed container presenting the connected groupoid $B(\Z/2)$ and a
right factor $\D\cong\Z/2$ of internal relabelings for which $[\omega]\neq0$; hence
no Zappa--Sz\'ep product / merge exists. Explicitly, take $\K=\Z/4$ and
$\D=\{0,2\}\cong\Z/2$. Then $\Sk_\C=\Z/4\,/\,\Z/2=\Z/2=B(\Z/2)$ is a connected
groupoid, \textnormal{(L)} and \textnormal{(H)} hold, and
\[
   [\omega]\;=\;\text{the generator of }\;H^2(\Z/2;\Z/2)\cong\Z/2 ,
\]
the class of the non-split extension $1\to\Z/2\to\Z/4\to\Z/2\to1$.
\end{theorem}
\begin{proof}
$\Z/4$ is abelian so $\D=\{0,2\}$ is abelian and normal; by
Propositions~\ref{prop:setup}--\ref{prop:schreier}, \textnormal{(L)} and
\textnormal{(H)} hold, $\Sk_\C=\Z/2$, the action of $\Z/2$ on $\D$ is trivial
(abelian $\K$), and $[\omega]$ is the extension class of $\Z/4$. The only order-$2$
subgroup of $\Z/4$ is $\D$ itself, so $\D$ has \emph{no} complement: the extension
is non-split and $[\omega]\neq0$. Since $H^2(\Z/2;\Z/2)\cong\Z/2$ (trivial action),
$[\omega]$ is its generator. All of this is machine-checked
(\S\ref{sec:comp}, example (B)): $\#\text{complements}=0$, the $\mathbb F_2$ bar
complex has $\dim H^2=1$, and $\omega_T\notin B^2$.
\end{proof}

\begin{remark}[Where the conjecture broke]\label{rem:broke}
The seed conjecture read ``connected groupoid $\Rightarrow\Sk_\C$ contractible
$\Rightarrow H^2=0$.'' A connected groupoid $\Sk_\C$ is equivalent, as a category,
to its vertex group $\Gamma$; its classifying space is $B\Gamma=K(\Gamma,1)$, which
is contractible \emph{iff $\Gamma=1$}. For $\Gamma=\Z/2$, $B\Gamma=\mathbb{RP}^\infty$
has $H^2(\mathbb{RP}^\infty;\Z/2)=\Z/2\neq0$. The broken assumption was the tacit
identification of ``every object isomorphic to every other'' (which \emph{does}
hold in a connected groupoid, making the \emph{skeleton} a single object) with
``the classifying space is a point'' (which needs the vertex group trivial too).
Invertibility collapses the objects but not the loops; the loops are precisely the
$H^{\ge1}$.
\end{remark}

\begin{theorem}[Same base, split total: obstruction is genuinely about $\K$]\label{thm:contrast}
The value of $[\omega]$ is not a function of the base groupoid alone. With the same
base $\Sk_\C=B(\Z/2)$ and the same coefficient $\D\cong\Z/2$, the total category
$\K=\Z/2\times\Z/2$ with $\D=\langle(1,0)\rangle$ has $[\omega]=0$: $\D$ has two
complements $\langle(0,1)\rangle,\langle(1,1)\rangle$, the split extension merges.
\end{theorem}
\begin{proof}
Machine-checked (\S\ref{sec:comp}, example (C)): $\#\text{complements}=2$,
$\omega_T\in B^2$, $[\omega]=0$. Analytically, $\Z/2\times\Z/2$ is the split
extension of $\Z/2$ by $\Z/2$, whose class is $0$.
\end{proof}

A larger, non-cyclic groupoid base obstructs in the same way.

\begin{theorem}[A rank-two groupoid base]\label{thm:q8}
Let $\K=Q_8$ (quaternion group) and $\D=Z(Q_8)=\{\pm1\}\cong\Z/2$. Then
$\Sk_\C=Q_8/\{\pm1\}=(\Z/2)^2$ is a connected groupoid, \textnormal{(H)} holds,
$H^2\big((\Z/2)^2;\Z/2\big)\cong(\Z/2)^3$, and $[\omega]\neq0$: every subgroup of
$Q_8$ contains $-1$, so $\D$ has no complement and the extension is non-split. No
merge exists.
\end{theorem}
\begin{proof}
Machine-checked (\S\ref{sec:comp}, example (D)): the $\mathbb F_2$ bar complex
returns $\dim H^2=3$ (so $|H^2|=8$, the known $H^*((\Z/2)^2;\Z/2)=\Z/2[x,y]$ with
$\dim H^2=3$), $\#\text{complements}=0$, and $\omega_T\notin B^2$. $Q_8$ is the
non-split central extension realising the class $x^2+xy+y^2$.
\end{proof}

\section{The correct positive theorem: cohomological dimension, not invertibility}
\label{sec:positive}

The conjecture's \emph{conclusion} --- ``these bases always compose'' --- is true
for the cases that matter, but for a reason invisible to the contractibility story.

\begin{theorem}[$\cd\le1$ forces the merge, for every $\D$]\label{thm:cd}
Suppose the skeleton $\Sk_\C$ has cohomological dimension $\le1$, i.e.\
$H^2(\Sk_\C;A)=0$ for every coefficient system $A:\Sk_\C^{\mathrm{op}}\to\Ab$.
Then for every right factor $\D$ satisfying \textnormal{(L)} and \textnormal{(H)},
$[\omega]=0$ and the Zappa--Sz\'ep product $\K=\C\bowtie\D$ exists.
In particular this holds whenever $\Sk_\C$ is a \emph{free groupoid} --- every
connected component has a free vertex group --- e.g.\ the free groupoid on any graph.
\end{theorem}
\begin{proof}
By \textnormal{(T3)} the class $[\omega]$ lives in $H^2(\Sk_\C;\D)$, which is $0$ by
hypothesis; then \textnormal{(G)} holds and \textnormal{(T2)} gives the product.
For the free-groupoid case: a connected groupoid is equivalent to its vertex group
$\Gamma$, and Baues--Wirsching cohomology is invariant under equivalence of
categories \cite{bw}, so $H^2(\Sk_\C;\D)\cong\prod_{\text{comp.}}H^2(\Gamma_i;M_i)$
(group cohomology, $M_i$ the stalk, a $\Gamma_i$-module by Lemma~\ref{lem:iso}). A
free group has cohomological dimension $\le1$ (Stallings--Swan
\cite{stallings,swan}: $\cd\Gamma\le1\iff\Gamma$ free), whence
$H^{\ge2}(\Gamma_i;M)=0$ for \emph{every} module $M$. Thus $H^2(\Sk_\C;\D)=0$.
\end{proof}

\begin{corollary}[Trichotomy for a connected groupoid base]\label{cor:trichotomy}
Let $\Sk_\C$ be a connected groupoid with vertex group $\Gamma$ and let $\D$ satisfy
\textnormal{(H)}, giving a $\Gamma$-module $M$ and an extension
$1\to M\to\K\to\Gamma\to1$. Then $[\omega]\in H^2(\Gamma;M)$ is its class, and:
\begin{enumerate}
\item[\textnormal{(a)}] \textbf{$\Gamma$ free} $(\cd\le1)$: $H^2(\Gamma;M)=0$ for
all $M$; the merge \emph{always} exists.
\item[\textnormal{(b)}] \textbf{$\Gamma$ trivial} $(\Gamma=1)$: $\Sk_\C$ is
genuinely contractible (the seed conjecture's true subcase); the merge exists. This
is (a) with $\Gamma$ free of rank $0$, and is realised by codiscrete groupoids
(\S\ref{sec:comp}, example (A)).
\item[\textnormal{(c)}] \textbf{$\Gamma$ has torsion}: $H^2(\Gamma;M)$ can be
nonzero (Theorems~\ref{thm:refute},\ref{thm:q8}); the merge exists iff the specific
extension splits.
\end{enumerate}
The dividing line is $\cd(\Gamma)$, not invertibility. Freeness --- the absence of
holonomy loops beyond dimension one --- is what trivialises the obstruction; a
non-free (torsion) symmetry, even though fully invertible, can carry a nonzero
degree-two class.
\end{corollary}

\begin{remark}[Unification with the acyclic examples]\label{rem:unify}
Both regimes are the \emph{same} invariant $H^2(\Sk_\C;\D)$; the base's homotopy
type chooses the mechanism. Over an \emph{acyclic} base (a poset/DAG, the free
category on an acyclic quiver) the nontrivial vertex group sits at a source, the
restriction maps vanish, and $H^2$ is cut out by branch-and-rejoin geometry --- the
$\Z/2$ rigid twist, the $\Z/n$ lot-cursor \cite{gobstr,orch,supply}. Over a
\emph{groupoid} base every restriction map is an isomorphism (Lemma~\ref{lem:iso})
and $H^2$ is group cohomology of $\Gamma$. The extremes meet in one slogan: the
merge is automatic exactly when $\Sk_\C$ is one-dimensional ($\cd\le1$) --- a
forest of handoffs, or a free groupoid of reversible ties --- and it can obstruct
the moment the handoff skeleton acquires a genuine degree-two cocycle, whether from
a rejoining acyclic branch or from a torsion loop.
\end{remark}

\section{Application: which social networks merge for free}\label{sec:social}

The result stages Neil's social-network exploration. Model a platform as a directed
container whose skeleton is the free (categorical or groupoidal) structure on its
tie-graph; merging two platforms sharing users is the Zappa--Sz\'ep composition
$\C\bowtie\D$ along the shared sub-structure $\D$ of relabelings/identifications.
The grading here is \textsc{computed}: ``a real network \emph{is} this
groupoid'' is a faithful abstraction, not a fidelity claim (the standing SEED-Q4
caveat).

\begin{corollary}[Mutual ties merge; the reason is dimension one]\label{cor:social}
A network of \emph{mutual} (symmetric) ties presents the \emph{free groupoid} on its
friendship graph. Its vertex groups are the fundamental groups of the graph's
connected components --- free groups (a graph is a $K(\text{free},1)$). By
Theorem~\ref{thm:cd}, $H^2=0$ for every coefficient system: \emph{mutual-tie
platforms always merge consistently}, for any right factor of shared
identifications. By contrast a network of \emph{directed} (follows) ties presents a
free \emph{category} on a directed quiver; when it is acyclic the earlier analysis
applies and a rejoining of two follow-paths can obstruct
($[\omega]=\varepsilon\in\Z/n$, \cite{orch,supply}). And a shared symmetry with
\emph{torsion} --- a finite cyclic identification cycle of order $n$ that is not
free --- can obstruct even in the mutual case (Theorem~\ref{thm:refute} is $n=2$).
\end{corollary}

\noindent The corrected headline for the grant's Impact line is therefore sharper
than the conjecture: \emph{reversibility alone does not guarantee that two platforms
merge; \textbf{freeness} does.} Mutual-tie graphs merge because their groupoids are
one-dimensional, not because they are trivial; introduce a torsion identification
(a quotient that forces a nontrivial finite loop) and the merge can fail with a
computable degree-two obstruction.

\section{Computational verification}\label{sec:comp}

All finite claims are machine-checked in \texttt{projects/scratch/groupoid-zs/}:
\texttt{groupoid\_zs.py} (orbits, \textnormal{(L)}, \textnormal{(H)}, $\Sk_\C$, the
defect, \#complements, and the $\mathbb F_2$ bar complex with the \emph{correct}
restriction maps $\varphi$), \texttt{dcont\_check.py} (D1--D5 for the bases). The
$\varphi$ computation matters: unlike the acyclic scripts, the groupoid restriction
maps are isomorphisms, not zero (Lemma~\ref{lem:iso}), so the earlier hard-coded
$\varphi=0$ is replaced by conjugation.

\begin{center}
\small
\begin{tabular}{llccccl}
\toprule
 & base $\Sk_\C$ & $\D$ & (L),(H) & $\dim H^2$ & \#compl. & verdict\\
\midrule
(A) codiscrete $\bullet\rightleftarrows\bullet$ & $1$ (contractible) & $\{\id\}$ & yes & $0$ & $1$ & merges\\
(B) $\Z/4$ & $B(\Z/2)$ & $\Z/2$ & yes & $1$ & $0$ & \textbf{obstructed}\\
(C) $\Z/2\times\Z/2$ & $B(\Z/2)$ & $\Z/2$ & yes & $1$ & $2$ & merges\\
(D) $Q_8$ & $B((\Z/2)^2)$ & $\Z/2$ & yes & $3$ & $0$ & \textbf{obstructed}\\
(E) $\Z/6$ & $B(\Z/2)$ & $\Z/3$ & yes & --- & $1$ & merges\\
\bottomrule
\end{tabular}
\end{center}

\noindent Examples (B),(D) are connected-groupoid bases with $[\omega]\neq0$
(Theorems~\ref{thm:refute},\ref{thm:q8}); (C) is the same base as (B) with a split
total (Theorem~\ref{thm:contrast}); (E) illustrates Schur--Zassenhaus (coprime
$|\D|,|\Gamma|$ force a split, $[\omega]=0$). In every case
$\#\text{complements}>0\iff[\omega]=0$, confirming Proposition~\ref{prop:schreier};
$\dim H^2$ reproduces the known group cohomologies $H^2(\Z/2;\Z/2)=\Z/2$,
$H^2((\Z/2)^2;\Z/2)=(\Z/2)^3$; and $\delta^2\delta^1=0$ is verified on the full
(now non-vacuous) $C^3$.

\section{Status, novelty, and honest scope}\label{sec:status}

\paragraph{What is proved.} (i) The reduction: over a one-object groupoid base
under \textnormal{(H)}, the merge obstruction is exactly the Schreier extension
class, and the merge exists iff the extension splits
(Propositions~\ref{prop:setup}--\ref{prop:schreier}). (ii) The refutation of the
strong conjecture: connected-groupoid bases $B(\Z/2)$ and $B((\Z/2)^2)$ with
$[\omega]\neq0$ (Theorems~\ref{thm:refute},\ref{thm:q8}). (iii) The correct positive
theorem: $\cd(\Sk_\C)\le1$ --- in particular a free groupoid --- forces the merge for
every $\D$ (Theorem~\ref{thm:cd}), giving the social-network corollary with the
right reason (Corollary~\ref{cor:social}). Grade: the worked \emph{instances} are
\textsc{computed} (machine-checked); the reduction and the $\cd\le1$ theorem are
\textsc{proved}, resting on the cited \textnormal{(T2)}, \textnormal{(T3)},
Schreier theory, and Stallings--Swan.

\paragraph{What is \emph{not} claimed (novelty firewall).} Group extensions and
their $H^2$ classification are classical (Schreier; \cite{maclane,brown}); the
identification $H^2_{\mathrm{BW}}(B\Gamma;M)=H^2(\Gamma;M)$ is standard. Groupoid
Zappa--Sz\'ep products themselves exist in the operator-algebra / self-similar
literature (Mundey--Sims and neighbours), on the \emph{completeness} axis, and are
cited as adjacent, \emph{not} a scoop. The contribution here is narrow and precise:
placing the groupoid case inside the Zappa--Sz\'ep merge program, \emph{refuting}
the seed contractibility conjecture, and isolating $\cd\le1$ (freeness) as the true
trivialising condition --- one honest correction, machine-witnessed, that redirects
the application story from ``reversible $\Rightarrow$ mergeable'' to ``free
$\Rightarrow$ mergeable.''

\paragraph{Open.} The nonabelian regime (\textnormal{(H)(ii)} fails: $\D$ non-normal,
i.e.\ a genuine two-sided Zappa--Sz\'ep matched pair rather than a semidirect
product) pushes the obstruction to a nonabelian $H^2$ of the matched pair
(Kac/Pirashvili, cited in \cite{gobstr}); we do not enter it. And the fidelity
question --- whether a given real network \emph{is} the free groupoid on its
tie-graph --- stays \textsc{open} (SEED-Q4).

\begin{thebibliography}{99}
\bibitem{pairwise} MacBeth, \emph{The pairwise Zappa--Sz\'ep criterion}, Kodamai
note, 10 June 2026.
\bibitem{gobstr} MacBeth, \emph{The global-closure obstruction (G) is a cohomology
class: the Zappa--Sz\'ep holonomy as $H^2(\Sk_\C;\D)$}, Kodamai note, 11 June 2026.
\bibitem{orch} MacBeth, \emph{The re-entrancy obstruction is the token-mutation bit},
Kodamai note, 20 July 2026.
\bibitem{supply} MacBeth, \emph{Supply-chain composition is a Zappa--Sz\'ep product},
Kodamai note, 23 July 2026.
\bibitem{rw} R.~Rosebrugh, R.~J.~Wood, \emph{Distributive laws and factorization},
J.\ Pure Appl.\ Algebra 175 (2002), 327--353.
\bibitem{bw} H.-J.~Baues, G.~Wirsching, \emph{Cohomology of small categories},
J.\ Pure Appl.\ Algebra 38 (1985), 187--211.
\bibitem{maclane} S.~Mac Lane, \emph{Homology}, Springer, 1963 (Ch.~IV: group
extensions and $H^2$).
\bibitem{brown} K.~S.~Brown, \emph{Cohomology of Groups}, GTM 87, Springer, 1982.
\bibitem{stallings} J.~Stallings, \emph{On torsion-free groups with infinitely many
ends}, Ann.\ of Math.\ 88 (1968), 312--334.
\bibitem{swan} R.~G.~Swan, \emph{Groups of cohomological dimension one},
J.\ Algebra 12 (1969), 585--610.
\end{thebibliography}

\end{document}
